\begin{python0}
from solutions import *; clear()
\end{python0}
\sectionthree{Structured vs OO}

We don't have classes and objects yet (what are they
anyway?) But even with what we have done so far, we can ask ...

Why? Why are we writing code this way.

The main reason is to control complexity ...

Structured programming focuses on functions. You start with a goal/task
(big function) and subdivide task into simpler subtask (simpler
functions). Data is passed between functions. Structured programming and structured
analysis tend to separate the computational task and the data involved.

OOP and OO analysis tend to focus on both data and functions together at
the same time. At a higher level of design and engineering, we think of
OO analysis and design as finding objects and their responsibilities
(functions), i.e., what the objects do. So in the discipline of OO
analysis and design, we tend to think of and find high level concepts
and analyze what the variable (i.e., the objects) associated with that
high level concept should be be capable of doing.

Let me give you an example.
Suppose I work for a bank and I need to print a daily report of all the
transactions (deposit to an account, withdrawal from an account, closing
of an account, opening of an account, changing the address of the
customer, etc.)

Structured thinking means this: I want to print a report. To do that, I need
to read the database for all transaction. For each transaction, there's an
transaction ID, the customer ID, a transaction type, etc. Now for each
customer ID, I will read the customer file, look for the matching customer
ID, read his firstname and lastname, etc. Next, for this transaction, I have
to look up the meaning of the transaction from the transaction type (for
instance say 1 means open account, 2 means close account, 3 means
deposit, etc.) I get the description of the transaction ID from some file, so
if the transaction is 1, the description for this entry is “Open Account”,
“Close Account”, etc. This style of thinking is called structured analysis
and design. And the type of code you write to reflect this way of thinking
tends to have a certain look – structured programming code.

The object-oriented way is different. You want to think of a report as
being made up of transaction objects. Each transaction object of course
has some type and some description (“Open Account”, “Close Account”,
etc.), but the main program will not get the transaction description.
Rather for each transaction, you will probably call

\begin{center}
\texttt{transaction.getDescription()}
\end{center}

in other words, each transaction object has the ability to provide useful
computations. You think of “telling the transaction to do its work on its 
own” and you wait for the description to be returned (probably a string).
Also, each transaction knows (on its own) that it's a transaction for a
particular customer. So to know more about the customer, you would ask
transaction to tell you who is the customer:
\begin{center}
\texttt{Customer customer = transaction.getCustomer();}
\end{center}

And with this customer object now available, you ask the customer to
give you his/her first and last name (probably string) say for printing:

\tab[3em]{\texttt{std::cout << customer.getLastName() << “, “}}
\tab[5em]{\texttt{<< customer.getFirstName()}}
\tab[5em]{\texttt{<< ...}}

So as you can see, the style of thinking (the OO analysis and design) will
give rise to very different code:

\tab[3em]{for each transaction in today's collection of transactions:}\\
 \tab[5em]{I ask the transaction for the transaction description and print it.}\\
 \tab[5em]{I then ask the transaction for the time when it was created and}\\
 \tab[5em]{print the date and time.}\\
 \tab[5em]{I then ask the transaction for firstname and lastname of the}\\
 \tab[5em]{customer responsible for this transaction this will result in the}\\
 \tab[5em]{transaction asking the cusomter involved for his/her firstname}\\
 \tab[5em]{and lastname. The customer gives his/her firstname and}\\
 \tab[5em]{lastname to the transaction and the transaction passes the}\\
 \tab[5em]{firstname and lastname back to me.}\\

Note that when I ask a transaction for the firstname and lastname of the
customer, the transaction talks to the customer (object) and asks the
customer to provide his/her first and last name. When an object needs to
perform some computation it will either do it himself/herself or will talk to
some other object to do the work.

You don't need to know the full picture of structured thinking and OO
thinking. The above is to give you a quick overview of the different
philosophy between the two.

The goal in this course is to teach you the basics of OO syntax and the
language features. You will NOT be able to engineer beautiful and well
constructed OO systems immediately. In fact it takes a very long time to
be an expert in either the structured or the object-oriented style of
analyzing and designing systems. Right now, my goal is to give you the
basic syntax and language features of an OO language. It will take many
years before you will become a true OO guru because the road is
challenging, tough, and is just extreme fun. (Which, frankly speaking, is a
good thing -- if OO is something you can learn in 1-2 years, you wouldn't
expect it to be fun or worth much, would you?)


